\documentclass[11pt]{article}

\usepackage[margin=1in]{geometry}
\usepackage{booktabs}
\usepackage{array}
\usepackage{tabularx}
\usepackage{hyperref}
\usepackage{xcolor}

\hypersetup{
  colorlinks=true,
  linkcolor=blue!50!black,
  urlcolor=blue!50!black,
  citecolor=blue!50!black
}

\newcommand{\code}[1]{\texttt{\small #1}}
\newcommand{\sysml}{SysML v2}
\emergencystretch=2em

\title{Executable-Rule Ablation on the Seeded GPT-5.5 Subset}
\author{SysML2-NL Ablation Study}
\date{}

\begin{document}
\maketitle

\begin{abstract}
We evaluate three \sysml{} generation sources on the same seeded 10\% subset
of the agent-generated corpus: the existing pipeline output (Ours),
Codex using GPT-5.5, and direct GPT-5.5 through the Vercel AI Gateway. The
evaluation uses the repository's text-first executable-rule checker,
\code{script/sysml\_executable\_rules.py}, over 155 samples selected with
seed 0 from \code{dataset/data/000387--001935}. Ours obtains the highest
aggregate executable-rule score, but all systems retain a substantial weakness
on signal-like accept-event outputs. The result suggests that the pipeline's
structural grounding improves executable-model compliance relative to direct
single-shot GPT-5.5 generation, while the remaining error mode is concentrated
enough to motivate targeted repair.
\end{abstract}

\section{Experimental Setup}

The subset is exactly the \code{--random\_10per} selection used by
\code{batch\_vercel\_gpt55.py}: 155 samples, drawn with \code{random.Random(0)}
from all eligible generated prompts. For each selected sample ID, the three
artifacts compared are:

\begin{center}
\small
\begin{tabularx}{\linewidth}{l>{\raggedright\arraybackslash}X>{\raggedright\arraybackslash}X}
\toprule
System & File pattern & Interpretation \\
\midrule
Ours & \code{dataset/data/<id>/<id>.sysml} &
Pipeline-generated corpus artifact \\
Codex+GPT-5.5 & \code{dataset/data/<id>/<id>.codex.sysml} &
Codex single-shot generation with GPT-5.5 \\
GPT-5.5 & \code{dataset/data/<id>/<id>.gpt55.sysml} &
Direct Vercel AI Gateway GPT-5.5 generation \\
\bottomrule
\end{tabularx}
\end{center}

Each artifact is scored by the local executable-rule checker,
\code{script/sysml\_executable\_rules.py}, on the five executable-modeling
rules shown in Table~2. The rule taxonomy follows the SysML 2 Rule
Verification Guide~\cite{sysml2RuleVerificationGuide}. The checker assigns
scores of 1.0 for \code{pass} and
\code{not\_applicable}, 0.5 for \code{unsupported}, and 0.0 for
\code{fail}. The aggregate score is the mean over all 155 samples and five
rules (\(775\) rule evaluations per system).

\section{Results}

\begin{table}[h]
\centering
\caption{Aggregate executable-rule outcomes on the selected 10\% subset.}
\begin{tabular}{lrrrrr}
\toprule
System & Mean score & All-pass samples & Pass & Fail & Unsupported \\
\midrule
Ours & 0.8968 & 80 / 155 & 695 & 80 & 0 \\
GPT-5.5 & 0.8710 & 57 / 155 & 675 & 100 & 0 \\
Codex+GPT-5.5 & 0.8600 & 47 / 155 & 666 & 108 & 1 \\
\bottomrule
\end{tabular}
\end{table}

Ours has the highest mean executable-rule score and the largest number of
samples passing all five checks. The direct GPT-5.5 run is slightly stronger
than the Codex+GPT-5.5 run in this measurement, but the margin is small.

\begin{table}[h]
\centering
\caption{Per-rule pass counts. Each row has 155 evaluations per system.}
\begin{tabular}{lrrr}
\toprule
Rule & Ours & GPT-5.5 & Codex+GPT-5.5 \\
\midrule
\code{ACCEPTEVENTOUTPUT} & 83 pass / 72 fail & 57 pass / 98 fail & 48 pass / 107 fail \\
\code{MESSAGEFLOWNEEDED} & 155 pass & 155 pass & 155 pass \\
\code{MESSAGESIGNATURE} & 154 pass / 1 fail & 155 pass & 155 pass \\
\code{STMINTEGRITY} & 148 pass / 7 fail & 153 pass / 2 fail & 154 pass / 1 fail \\
\code{SUBMACHINESTR} & 155 pass & 155 pass & 154 pass / 1 unsupported \\
\bottomrule
\end{tabular}
\end{table}

The dominant failure mode is \code{ACCEPTEVENTOUTPUT}. Ours reduces this
failure count from 98--107 in the GPT-5.5 baselines to 72, but it does not
eliminate the issue. Conversely, the GPT-5.5 baselines are slightly cleaner on
\code{STMINTEGRITY}, where Ours has seven failures. This indicates that the
pipeline improves event-payload modeling more than it improves state-machine
call ownership.

\begin{table}[h]
\centering
\caption{Paired per-sample score differences. Confidence intervals are
descriptive normal approximations over the 155 paired samples.}
\begin{tabular}{lrrr}
\toprule
Comparison & Mean difference & 95\% CI & Better / equal / worse \\
\midrule
Ours $-$ Codex+GPT-5.5 & +0.0368 & [0.0128, 0.0607] & 60 / 66 / 29 \\
Ours $-$ GPT-5.5 & +0.0258 & [0.0026, 0.0490] & 49 / 78 / 28 \\
GPT-5.5 $-$ Codex+GPT-5.5 & +0.0110 & [-0.0105, 0.0325] & 40 / 85 / 30 \\
\bottomrule
\end{tabular}
\end{table}

The paired analysis supports two conclusions. First, Ours has a consistent but
modest advantage over both single-shot baselines on executable-rule compliance.
Second, the difference between direct GPT-5.5 and Codex+GPT-5.5 is not robust
under this metric: most samples are tied, and the confidence interval spans
zero.

\section{Interpretation}

The ablation points to a concentrated compliance gap rather than broad
syntactic degradation. Message-flow realization, message signatures, and
submachine structure are almost saturated for all three systems. The main
remaining issue is whether signal-like \code{accept} constructs expose a typed
output or payload. This is a semantically narrow but executable-critical
failure: a model can look structurally complete while still omitting the data
interface needed by downstream behavior.

The pipeline output is strongest overall, plausibly because its generation path
uses retrieval, synthesis, and syntax-oriented filtering rather than a single
completion. However, its lower \code{STMINTEGRITY} score shows that this
advantage is not uniform. A practical next repair pass should therefore be
rule-targeted: add a post-generation check that rewrites signal-like
\code{accept} statements with explicit typed payloads, while separately
verifying that state-machine entry, exit, effect, and \code{perform} calls
resolve to available action or operation definitions.

\section{Limitations}

The executable-rule checker is deliberately text-first and approximate; it is
useful for bulk comparison but is not a substitute for a full \sysml{} parser
or semantic resolver. The subset is fixed by seed 0 and contains only generated
corpus prompts, so the results should be interpreted as an ablation on this
study slice rather than a claim about all SysML modeling tasks. Finally, the
scoring function treats \code{not\_applicable} as pass and \code{unsupported}
as half credit, which is appropriate for screening but can hide differences in
model expressiveness.

\bibliographystyle{plain}
\bibliography{references}

\end{document}
